\section{Bounds on claims and conditions for deployment}\label{sec:C30:07}
Repeated improvement creates two kinds of accumulated state: executable changes and beliefs about why those changes help. Both need rollback. An acceptance record should include the hypothesis, modified artifacts, compatibility tuple, validation outcomes, held-out outcomes, resource cost, known regressions, and the exact conditions for reversal. Returning to old code while retaining a newly trained model or incompatible memory does not restore the old system.

Deployment should be gated on the property being changed. A context-saving edit needs tests for rare decision-critical information and conditioning consistency. A skill update needs applicability and permission checks. A learned router needs no-valid-tool and schema-drift tests. A verifier update needs independent adversarial cases. A model update needs behavioral and memory-interpretation regression tests. A global score can supplement these checks but cannot replace them, because many failures occupy a small probability mass with high consequence.

The strongest current evidence supports bounded adaptation mechanisms under specified task and feedback conditions. The harder open problem is maintaining trustworthy selection as executor capability, evaluator behavior, state representations, and environments co-evolve. A system that becomes increasingly effective at satisfying a fixed evaluator may be improving the intended task, exploiting the evaluator, or both. The decisive evidence comes from independent transfer, mechanism ablations, preserved failure records, and explicit evaluation of the improvement process itself.
